% !TeX root = ../../Book2.tex
% 纲要标题：## 第24章*　拓扑 K 理论、Bott 周期性与指标的拓扑侧
% 纲要位置：几何学计划纲要.md，第 1923 行。
% 正文状态：初稿；尚未正式审核。

% 纲要内容开始
% 
% > 本章目标与证明范围：以有限 CW 空间上的复 K 理论为完整主模块，基本乘积定理的证明可在紧 Hausdorff 参数空间上进行。Bott 采用多项式粘接的具体证明路线；实 KO 的八周期和一般指标定理作为准确外部输入，不以复数特例混称全理论。
% 
% 纲要内容结束

\OptionalChapter{拓扑 K 理论、Bott 周期性与指标的拓扑侧}{g2:ch:24}
\BookChapterNumber{g2:ch:24}

\begin{chapterabstract}
	向量丛的直和可以做加法，却通常不能直接做减法。群完备化把它变成虚丛的群，张量积进一步给出环结构。本章建立相对群、悬挂和长正合列，并用多项式粘接完整证明复 Bott 周期性：先把圆周上的连续粘接函数逼近为 Laurent 多项式，再把高次矩阵线性化，最后由单位圆内外的谱子丛得到基本乘积定理及其逆映射。陈特征把有限 CW 空间的有理 K 群与偶奇上同调联系起来；椭圆主符号则给出余切丛上的紧支 K 类，通向解析指标和拓扑指标的比较。
\end{chapterabstract}

\begin{learninggoals}
	\begin{itemize}
		\item 从向量丛的直和构造 \(K^0\)，区分实际同构、稳定同构与虚丛相等，并计算低维球面的例子。
		\item 用映射锥和三元组理解相对 K 群，构造悬挂、连接映射及长正合列。
		\item 完成粘接函数的多项式逼近、线性化及连续谱分解，构造 \(K^0(X\times S^2)\) 分解的两个方向。
		\item 推出复 K 理论的二周期、球面系数及有限 CW 算例，辨认实 KO 八周期使用的不同输入。
		\item 证明陈特征的乘法性与有理同构，理解椭圆符号、紧支类和指标推前各自记录的信息。
	\end{itemize}
\end{learninggoals}

第19章把球面上的向量丛描述为沿赤道粘接两块平凡丛，第21章又用陈类测量复丛的扭曲。现在考虑一个更具体的问题：\(S^2\) 上的 Hopf 线丛 \(H\) 与平凡线丛有什么“差”？两者的秩相同，却有不同的第一陈类。若能把 \([H]-[\underline{\mathbb C}]\) 当作一个可运算的对象，就能把秩与扭曲分开，并在直和、张量积和拉回中追踪这份差异。这正是 K 理论引入虚丛的出发点。

\ProseParagraphBreak

更令人惊讶的是，把任意紧参数空间 \(X\) 乘上一张二球面，并不会产生难以控制的全新类型。令 \(\beta=[H]-1\)，本章将证明每个虚丛都能唯一写成 \(a+b\beta\)，其中 \(a,b\in K^0(X)\)。\(a\) 是在球面一点上的限制，\(b\) 则保留粘接函数绕单位圆变化时的整族信息。对单个矩阵而言，单位圆内特征值的个数是一个整数；对连续矩阵族而言，相应广义特征空间拼成一条向量丛。这一步把线性代数与全局拓扑联系起来，也是证明中最值得看清的构造。

\ProseParagraphBreak

主要准备是第19章的丛粘接、同伦不变性和平凡丛补空间，第1章的映射锥及同伦延拓性质，以及第21章的陈类和分裂原理。长证明中用到的复矩阵谱投影会在原处由 Jordan 块和圆周积分展开。最后一节把本章所得用于椭圆算子的主符号；解析 Fredholm 性及一般指标等式在使用处给出条件和正式文献，实 KO 的八周期也在相应陈述后注明来源。

% !TeX root = ../../Book2.tex
% 纲要标题：### 24.1 向量丛的群完备化
% 纲要位置：几何学计划纲要.md，第 1927 行。
% 正文状态：初稿；尚未正式审核。

% 纲要内容开始
% 
% * K⁰、约化 K⁰ 与相对群
% * 稳定同构和虚丛
% * 粘合构造与球面 K 群
% * 环结构和外积
% 
% 纲要内容结束

\section{向量丛的群完备化}
\label{g2:sec:2401}

第19章把向量丛分类为到 Grassmann 空间的映射，却还没有把全部丛组织成一个可作减法的对象。直和是现成的加法，但一张正秩丛不可能通过直和得到零丛，因此通常没有逆元。本节把“比较两张丛的差”作为新的对象：增加同一张辅助丛不会改变差，稳定同构由此成为等式，张量积则成为乘法。后面正合列与周期性所计算的，正是这套稳定的丛信息。

\subsection{\texorpdfstring{\(K^0\)、约化 \(K^0\) 与相对群}{K⁰、约化 K⁰ 与相对群}}
\label{g2:subsec:2401:01}
本节的底空间先取紧 Hausdorff 空间 \(X\)，丛取有限秩连续复向量丛。允许秩随连通分支变化；秩函数局部常值，紧性保证它只取有限多个值。记 \(\operatorname{Vect}_{\mathbb C}(X)\) 为这些丛的同构类所成集合，包含零丛。直和使它成为交换幺半群：运算满足结合律、交换律并有零元，但尚未保证逆元，也不能预先作消去。

若把 \((E,F)\) 理解为形式差 \(E-F\)，那么 \((E\oplus G,F\oplus G)\) 应表示同一个差。不过只允许这种直接增添，还不足以让所有形式差自然组成群；下面的关系把所需稳定比较一次写全。

\begin{definition}\label{g2:def:topological-k-zero}
	对紧 Hausdorff 空间 \(X\)，在复向量丛对上规定
	\[
		(E,F)\sim(E',F')
		\quad\Longleftrightarrow\quad
		E\oplus F'\oplus G\cong E'\oplus F\oplus G
	\]
	对某张复向量丛 \(G\to X\) 成立。所得等价类集合记作 \(K^0(X)\)，称为 \(X\) 的复\term[tuopu-K-qun]{拓扑 K 群}[topological K-group]。丛对的类写作 \([E]-[F]\)，其中 \([E]\) 表示 \((E,0)\) 的类。运算为
	\[
		([E]-[F])+([E']-[F'])=[E\oplus E']-[F\oplus F'].
	\]
	这样的形式差称为\term[xucong]{虚丛}[virtual bundle]。\symidx{\(K^0(X)\)：复向量丛的群完备化}
\end{definition}

\begin{proposition}\label{g2:prop:k-group-completion}
	上述关系为等价关系，\(K^0(X)\) 是交换群，其零元为 \([0]\)，\([E]-[F]\) 的逆元为 \([F]-[E]\)。映射
	\[
		j:\operatorname{Vect}_{\mathbb C}(X)\longrightarrow K^0(X),\qquad E\longmapsto[E]
	\]
	具有如下泛性质：对任意阿贝尔群 \(A\) 及保持直和、零元的映射 \(a:\operatorname{Vect}_{\mathbb C}(X)\to A\)，存在唯一群同态 \(\overline a:K^0(X)\to A\) 满足 \(\overline a\circ j=a\)。
\end{proposition}
\begin{proof}
	自反性可取 \(G=0\)，对称性由交换等式两侧得到。若比较 \((E,F)\)、\((E',F')\) 使用辅助丛 \(G\)，比较 \((E',F')\)、\((E'',F'')\) 使用 \(G'\)，则先后代入两份同构可得
	\[
		E\oplus F''\oplus(F'\oplus G\oplus G')
		\cong E''\oplus F\oplus(F'\oplus G\oplus G').
	\]
	这证明传递性，并显示为什么关系中保留辅助丛。对两份稳定同构取直和，就得到它们之和的稳定同构，所以加法良定义。结合与交换来自丛直和的相应同构，而 \((E\oplus F,F\oplus E)\sim(0,0)\)，故所列逆元成立。

	定义 \(\overline a([E]-[F])=a(E)-a(F)\)。若两对等价，对定义中的同构应用 \(a\)，在群 \(A\) 中消去 \(a(G)\)，便得到相同的差。因此它良定义且可加。每个元素都是两个 \(j\) 的值之差，所以这一要求也确定了唯一的 \(\overline a\)。相应交换图为
	\[
		\begin{tikzcd}[column sep=large]
			\operatorname{Vect}_{\mathbb C}(X)\arrow[r,"j"]\arrow[dr,"a"']&K^0(X)\arrow[d,dashed,"\exists!\,\overline a"]\\
			&A.
		\end{tikzcd}
	\]
\end{proof}

这个构造称为幺半群的\term[qun-wanbeihua]{群完备化}[group completion]。它把所有可加的丛不变量统一到同一个群上。例如，单点上的复向量空间只由维数决定，故 \(K^0(\{*\})\cong\mathbb Z\)。对一般 \(X\)，逐点取秩给出
\[
	\operatorname{rk}:K^0(X)\longrightarrow C(X,\mathbb Z),\qquad
	\operatorname{rk}([E]-[F])(x)=\dim E_x-\dim F_x,
\]
其中 \(C(X,\mathbb Z)\) 是连续整数值函数群。这个整数差称为虚秩。连通时它是单个整数；不连通时各分支可以不同。

\begin{definition}\label{g2:def:reduced-relative-k}
	对带基点的紧 Hausdorff 空间 \((X,x_0)\)，定义\term[yuehua-K-qun]{约化 K 群}[reduced K-group]
	\[
		\widetilde K^0(X)=\ker\bigl(i^*:K^0(X)\longrightarrow K^0(\{x_0\})=\mathbb Z\bigr),
	\]
	其中 \(i\) 为基点包含，\(i^*\) 取纤维的虚维数。对有限 CW 对 \((X,A)\)，定义\term[xiangdui-K-qun]{相对 K 群}[relative K-group]
	\[
		K^0(X,A)=\widetilde K^0(X/A).
	\]
	若 \(A\ne\varnothing\)，商空间以压缩后的 \(A\) 为基点；若 \(A=\varnothing\)，约定此处的 \(X/A\) 为 \(X_+=X\sqcup\{*\}\)。\symidx{\(\widetilde K^0(X)\)：基点纤维虚秩为零的 K 群}\symidx{\(K^0(X,A)\)：有限 CW 对的相对 K 群}
\end{definition}

记 \(\varepsilon^r=X\times\mathbb C^r\)，并以 \(1=[\varepsilon^1]\) 表示平凡线丛的类。基点取秩有分裂 \(m\mapsto m1\)，故
\begin{equation}\label{g2:eq:k-rank-splitting}
	K^0(X)\cong\mathbb Z\oplus\widetilde K^0(X),\qquad
	x\longmapsto\bigl(\operatorname{rk}_{x_0}x,\ x-(\operatorname{rk}_{x_0}x)1\bigr).
\end{equation}
当 \(X\) 不连通时，约化只消去所选基点处的虚秩，其他分支上的秩仍可非零。例如 \(K^0(S^0)\cong\mathbb Z^2\)，选其中一点为基点后有 \(\widetilde K^0(S^0)\cong\mathbb Z\)。又因为 \(K^0(X_+)\cong K^0(X)\oplus\mathbb Z\)，空子空间约定给出 \(K^0(X,\varnothing)=K^0(X)\)，而 \(K^0(X,X)=0\)。下一节将把商空间定义改写为 \(A\) 上带有指定同构的数据，从而构造相对映射和连接同态。

\subsection{稳定同构和虚丛}
\label{g2:subsec:2401:02}
形式差的两个成员可以简化为“一张实际丛减去一个平凡丛”。其几何依据是第19章的补丛构造：由\cref{g2:prop:bundle-embedding-classifier}，紧基底上的有限秩丛可嵌入某个 \(\varepsilon^N\)，正交补就是所需补丛。局部常秩情形按有限多个开闭秩层分别构造，再把环境维数增到同一个 \(N\)，所以同样成立。

\begin{proposition}\label{g2:prop:k-stable-equality}
	设 \(X\) 紧且 Hausdorff。每个 \(x\in K^0(X)\) 都可写成 \([E]-[\varepsilon^N]\)。对两张复向量丛 \(E,F\to X\)，有
	\[
		[E]=[F]\text{ 于 }K^0(X)
		\quad\Longleftrightarrow\quad
		E\oplus\varepsilon^N\cong F\oplus\varepsilon^N
	\]
	对某个 \(N\ge0\) 成立。因此群中的实际丛等式恰好是第19章所说的稳定同构。
\end{proposition}
\begin{proof}
	给定 \([E]-[F]\)，取补丛 \(F'\) 使 \(F\oplus F'\cong\varepsilon^N\)。在形式差的两项中加上 \(F'\)，便得
	\[
		[E]-[F]=[E\oplus F']-[\varepsilon^N].
	\]
	若 \([E]=[F]\)，定义给出某个 \(G\) 使 \(E\oplus G\cong F\oplus G\)。再取 \(G'\) 使 \(G\oplus G'\cong\varepsilon^N\)，在两边加 \(G'\) 即得稳定同构。反过来，存在这样的稳定同构时，在等价关系中直接取 \(G=\varepsilon^N\) 即可。
\end{proof}

特别地，在带基点的空间上，每个 \(\widetilde K^0(X)\) 的类可写成 \([E]-[\varepsilon^N]\)，且 \(\operatorname{rank}E_{x_0}=N\)。若 \(X\) 连通，\(E\) 在所有点的秩都为 \(N\)。对连通 \(X\)，两项 \([E]-\operatorname{rank}E\) 与 \([F]-\operatorname{rank}F\) 相同，正是说两丛补加适当数量的平凡线以后同构；补加数量可不同，是因为已经减掉各自的秩。

\ProseParagraphBreak

虚丛的“负号”来自群运算。例如 \(-[E]=[0]-[E]\) 的虚秩是负数，而对偶丛 \(E^*\) 与 \(E\) 的实际秩相同。因而对偶通常不是直和的逆元。另一方面，若一个短正合列
\[
	0\longrightarrow E'\longrightarrow E\longrightarrow E''\longrightarrow0
\]
由 \(X\) 上有限秩复向量丛组成，那么 Hermitian 度量使 \(E'\) 的正交补映到 \(E''\) 的每根纤维上成为同构，故 \(E\cong E'\oplus E''\)，从而 \([E]=[E']+[E'']\)。这里讨论连续复向量丛；所取正交补是连续的，论证并没有要求它保持额外的全纯结构。

\begin{proposition}\label{g2:prop:k-pullback-homotopy}
	连续映射 \(f:X\to Y\) 在紧 Hausdorff 空间之间诱导群同态
	\[
		f^*:K^0(Y)\longrightarrow K^0(X),\qquad
		f^*([E]-[F])=[f^*E]-[f^*F].
	\]
	它满足 \((g\circ f)^*=f^*g^*\)、\(\operatorname{id}^*=\operatorname{id}\)，同伦的映射诱导相同同态。带基点映射还保持约化群，有限 CW 对的映射通过商空间诱导相对群同态。
\end{proposition}
\begin{proof}
	拉回保持丛同构与直和，所以保持定义中的等价关系和群运算。复合拉回与逐次拉回之间的纤维识别给出函子等式。若 \(f_0\simeq f_1\)，把任一 \(Y\) 上的丛沿同伦拉回到 \(X\times I\)，\cref{g2:thm:bundle-homotopy-invariance}说明两端限制同构，故它们在 \(K^0\) 中的类相同。基点纤维与商空间映射的自然性分别给出最后两项。
\end{proof}

可缩的非空紧 Hausdorff 空间于是有 \(K^0(X)\cong\mathbb Z\)，且任一基点下 \(\widetilde K^0(X)=0\)。具体地，收缩把每张丛同构到一个常值纤维丛；从此例开始，丛的稳定计算与空间的同伦类型联系起来。

\subsection{粘合构造与球面 K 群}
\label{g2:subsec:2401:03}
球面可分为两个可缩半球，丛在各半球都能平凡化，所以差异只留在赤道的粘合函数中。\cref{g2:thm:clutching-classification}已经证明，对 \(m\ge2\)，秩 \(r\) 复丛的同构类由
\[
	\pi_{m-1}(\operatorname{U}(r))
\]
分类。这里先用极分解把 \(\operatorname{GL}_r(\mathbb C)\) 缩到 \(\operatorname{U}(r)\)，再把赤道基点处的矩阵乘回单位矩阵，把自由同伦化为基点同伦。若 \(g\) 是粘合函数，则加一条平凡线把它变为 \(\operatorname{diag}(g,1)\)。

\begin{proposition}\label{g2:prop:k-clutching-spheres}
	对 \(m\ge2\)，粘合构造给出自然群同构
	\[
		\varinjlim_r\pi_{m-1}(\operatorname{U}(r))
		\xrightarrow{\ \cong\ }\widetilde K^0(S^m),
		\qquad[g]\longmapsto[E_g]-r.
	\]
	直接极限沿 \(g\mapsto\operatorname{diag}(g,1)\) 取；两项代表相同，是指它们在某个共同的有限秩阶段同伦。
\end{proposition}
\begin{proof}
	上一命题把约化类写成 \([E]-r\)，球面连通保证 \(\operatorname{rank}E=r\)，故粘合分类给出满射。若两粘合丛给出相同约化类，\cref{g2:prop:k-stable-equality}说明补加平凡线后两丛同构。再用固定秩的粘合分类，就得到稳定后的粘合函数同伦。这证明单射。

	直和对应分块矩阵 \(\operatorname{diag}(g,h)\)。把两块先稳定到同阶，令
	\[
		R_t=\begin{pmatrix}\cos t\,I&-\sin t\,I\\\sin t\,I&\cos t\,I\end{pmatrix},\qquad
		0\le t\le\pi/2.
	\]
	则 \(\operatorname{diag}(g,I)R_t\operatorname{diag}(I,h)R_t^{-1}\) 从 \(\operatorname{diag}(g,h)\) 连到 \(\operatorname{diag}(gh,I)\)，且在赤道基点始终为单位矩阵。拓扑群的球面映射中，逐点乘法与同伦群的拼接相同：把代表写在立方体上，两种运算具有同一单位元并满足逐坐标的交换律
	\[
		(f_1*f_2)(g_1*g_2)=(f_1g_1)*(f_2g_2),
	\]
	其中 \(*\) 表示沿同一个坐标拼接。将某两项取常值单位映射，就在同伦类上得到逐点乘积等于拼接。因此分块和表示同伦群的和，所构造双射为群同构。
\end{proof}

低维计算已经可以完成。第19章说明圆周上的复丛全部平凡，故 \(K^0(S^1)\cong\mathbb Z\)、\(\widetilde K^0(S^1)=0\)。对 \(S^2\)，\cref{g2:prop:classical-low-homotopy}给出
\[
	\det_*:\pi_1(\operatorname{U}(r))\xrightarrow{\ \cong\ }\pi_1(S^1)=\mathbb Z,
\]
且增加单位矩阵块保持这个同构。因此 \(\widetilde K^0(S^2)\cong\mathbb Z\)，分类整数就是 \(\det g\) 的绕数。由\cref{g2:prop:first-classes-determinants,g2:prop:chern-transition-comparison}，它也等于 \(\langle c_1(E_g),[S^2]\rangle\)，只须保持同一组赤道方向与坐标约定。

下面固定与第21章一致的正生成元。将 \(S^2\) 识别为复定向的 \(\mathbb CP^1\)，取
\[
	z=z_1/z_0,\qquad w=z_0/z_1=1/z,\qquad H=\gamma^*.
\]
对典范线丛 \(\gamma\)，局部标架为 \(e_z=(1,z)\)、\(e_w=(w,1)=z^{-1}e_z\)；对偶标架因而满足 \(e_w^*=z e_z^*\)。沿用第19章的关系 \((x,v)_+\sim(x,g(x)v)_-\)，取 \(D_+\) 为 \(w\) 圆盘、\(D_-\) 为 \(z\) 圆盘，并以 \(z=e^{\mathrm i\theta}\)、\(\theta\) 增大参数化赤道，则 \(g(z)=z\) 的粘合丛正是 \(H\)。第21章已经核定 \(c_1(H)=h\)、\(\langle h,[\mathbb CP^1]\rangle=1\)，所以
\begin{equation}\label{g2:eq:k-two-sphere-additive}
	\widetilde K^0(S^2)=\mathbb Z\,([H]-1),\qquad
	K^0(S^2)=\mathbb Z\,1\oplus\mathbb Z\,([H]-1).
\end{equation}
秩 \(r\ge1\)、第一陈数为 \(d\) 的丛满足 \([E]=r+d([H]-1)\)。事实上它还同构于 \(H^{\otimes d}\oplus\varepsilon^{r-1}\)：这张丛的粘合矩阵 \(\operatorname{diag}(z^d,1,\ldots,1)\) 与 \(g\) 有相同行列式绕数，故在 \(\operatorname{U}(r)\) 中同伦。负数 \(d\) 按对偶张量幂理解。

\begin{example}\label{g2:exm:k-sphere-tangent-line}
	\(\mathbb CP^1\) 的复切线丛满足 \(c_1(\mathrm T\mathbb CP^1)=2h\)，见式\eqref{g2:eq:complex-projective-tangent-classes}，因此
	\[
		[\mathrm T\mathbb CP^1]-1=2([H]-1)\ne0.
	\]
	球面的实切丛加一条实平凡线可平凡化，与这个结论并不冲突：这里用的是复一维切丛，稳定操作只允许加入复平凡方向。实、复理论的稳定信息不同，后面讨论周期性时也须保留基域。
\end{example}

\subsection{环结构和外积}
\label{g2:subsec:2401:04}
直和给出加法，张量积给出另一种运算。它在实际丛上满足 \((E\oplus F)\otimes G\cong(E\otimes G)\oplus(F\otimes G)\)，所以可以把张量积按分配律延到虚丛。

\begin{proposition}\label{g2:prop:k-ring-structure}
	对紧 Hausdorff 空间 \(X\)，下式把 \(K^0(X)\) 变为含幺交换环：
	\[
		([E]-[F])([G]-[J])
		=[E\otimes G\oplus F\otimes J]-[E\otimes J\oplus F\otimes G].
	\]
	单位元为 \(1=[\varepsilon^1]\)。连续映射的拉回是含幺环同态；基点取秩也是含幺环同态，因而 \(\widetilde K^0(X)\) 是 \(K^0(X)\) 的理想。
\end{proposition}
\begin{proof}
	固定一张实际丛 \(G\)，映射 \(E\mapsto[E\otimes G]\) 对直和可加，故由群完备化的泛性质唯一延拓为 \(K^0(X)\to K^0(X)\) 的同态。该同态又对 \(G\) 的直和可加，因此对第二个变量再次用泛性质，得到唯一的双线性乘法。展开两份差，就得到所列公式，也同时证明它与代表选择无关。

	实际丛的张量积有结合、交换及单位同构，这些等式对形式差双线性延拓，故得到环公理。拉回保持直和、张量积及平凡线丛，因而保持乘法和单位。纤维维数满足 \(\dim(E_x\otimes G_x)=\dim E_x\dim G_x\)，所以取秩也乘法相容，其核便为理想。
\end{proof}

\begin{definition}\label{g2:def:k-external-product}
	设 \(X,Y\) 为紧 Hausdorff 空间，\(p_X,p_Y\) 为积空间的投影。对 \(x\in K^0(X)\)、\(y\in K^0(Y)\)，定义\term[K-lilun-waiji]{K 理论外积}[external product in K-theory]
	\[
		x\boxtimes y=p_X^*x\,p_Y^*y\in K^0(X\times Y).
	\]
	对实际丛，它由 \(E\boxtimes F=p_X^*E\otimes p_Y^*F\) 表示。\symidx{\(x\boxtimes y\)：K 理论外积}
\end{definition}

外积可加于每个变量，且对连续映射 \(f:X'\to X\)、\(g:Y'\to Y\)，有
\begin{equation}\label{g2:eq:k-external-naturality}
	(f\times g)^*(x\boxtimes y)=f^*x\boxtimes g^*y.
\end{equation}
这是因为 \(p_X\circ(f\times g)=f\circ p_{X'}\)，另一投影亦然，再使用拉回的乘法相容性。若 \(X=Y\)，沿对角映射 \(\Delta:X\to X\times X\) 拉回便得到原来的乘法：\(\Delta^*(x\boxtimes y)=xy\)。交换两因子的映射则把 \(x\boxtimes y\) 送到 \(y\boxtimes x\)；当前两项均为零次，尚无分次符号。

\ProseParagraphBreak

若 \(X,Y\) 有基点，\(x\in\widetilde K^0(X)\)、\(y\in\widetilde K^0(Y)\)，那么 \(x\boxtimes y\) 在 \(X\times\{y_0\}\) 与 \(\{x_0\}\times Y\) 上的限制均为零。这为把外积看作压缩积 \(X\wedge Y\) 上的类提供了必要条件；下一节将通过相对正合性构造相应的约化外积，并证明其与这里的绝对外积相容。

在 \(S^2\) 的例子中，我们已得到 \(1\) 与 \([H]-1\) 组成的加法基。环结构还要求计算 \(([H]-1)^2\)，而积空间 \(X\times S^2\) 上的丛更需要知道两项系数是否都能从 \(K^0(X)\) 取得。这两件事将在\secref{g2:sec:2403}中由多项式粘合构造一起解决。
% 纲要条目排为编号小节标题，并进入目录与书签；说明段仍保留于纲要注释中。

% !TeX root = ../../Book2.tex
% 纲要标题：### 24.2 正合性和悬挂
% 纲要位置：几何学计划纲要.md，第 1934 行。
% 正文状态：未撰写；本文件目前是框架。

% 纲要内容开始
% 
% * 映射锥与相对 K 群
% * 悬挂定义负次数
% * 长正合列和切除
% * 乘积结构的自然性
% 
% 纲要内容结束

\section{正合性和悬挂}
\label{g2:sec:2402}

待撰写

\subsection{映射锥与相对 K 群}
\label{g2:subsec:2402:01}
待撰写

\subsection{悬挂定义负次数}
\label{g2:subsec:2402:02}
待撰写

\subsection{长正合列和切除}
\label{g2:subsec:2402:03}
待撰写

\subsection{乘积结构的自然性}
\label{g2:subsec:2402:04}
待撰写
% 纲要条目排为编号小节标题，并进入目录与书签；说明段仍保留于纲要注释中。

% !TeX root = ../../Book2.tex
% 纲要标题：### 24.3 复 Bott 周期性的多项式粘接证明
% 纲要位置：几何学计划纲要.md，第 1941 行。
% 正文状态：未撰写；本文件目前是框架。

% 纲要内容开始
% 
% 主模块固定有限 CW 空间上的复 K 理论，采用多项式粘接—基本乘积定理的证明路线；相对版本和胞腔粘合逐项说明。[R8]
% 
% * $\mathbb{CP}^1$ 上的 Hopf 线丛与 Bott 元
% 
%   取 $\beta=[H]-1\in\widetilde K^0(S^2)$，其中 $H$ 为所选 Hopf 线丛，并核验 $(H-1)^2=0$。
% 
% * $X\times S^2$ 上向量丛的粘接函数模型
% * Laurent 多项式逼近与粘接同伦的保持
% * 增添平凡因子后的多项式粘接线性化
% * 单位圆内外的谱分解与连续子丛
% * $K^0(X\times S^2)$ 的基本乘积定理及逆映射
% 
%   构造并证明 $K^0(X)\otimes_{\mathbb Z}\mathbb Z[H]/(H-1)^2\longrightarrow K^0(X\times S^2)$ 为同构。满射由多项式粘接约化得到，单射由逆映射及其同伦无关性得到；不能用未证的 K 理论 Künneth 公式替代。
% 
% * 相对群、约化外积与二次悬挂同构
% 
%   由积空间到楔和及压缩积的相对正合列推出 $\beta\boxtimes-:\widetilde K^0(X)\xrightarrow{\sim}\widetilde K^0(\Sigma^2X)$，并核验所用自然性。
% 
% * 复周期性、球面系数及有限 CW 空间的计算
% 
% 纲要内容结束

\section{复 Bott 周期性的多项式粘接证明}
\label{g2:sec:2403}

待撰写

\subsection{\texorpdfstring{\(\mathbb{CP}^1\) 上的 Hopf 线丛与 Bott 元}{CP¹ 上的 Hopf 线丛与 Bott 元}}
\label{g2:subsec:2403:01}
待撰写

\subsection{\texorpdfstring{\(X\times S^2\) 上向量丛的粘接函数模型}{X×S² 上向量丛的粘接函数模型}}
\label{g2:subsec:2403:02}
待撰写

\subsection{Laurent 多项式逼近与粘接同伦的保持}
\label{g2:subsec:2403:03}
待撰写

\subsection{增添平凡因子后的多项式粘接线性化}
\label{g2:subsec:2403:04}
待撰写

\subsection{单位圆内外的谱分解与连续子丛}
\label{g2:subsec:2403:05}
待撰写

\subsection{\texorpdfstring{\(K^0(X\times S^2)\) 的基本乘积定理及逆映射}{K⁰(X×S²) 的基本乘积定理及逆映射}}
\label{g2:subsec:2403:06}
待撰写

\subsection{相对群、约化外积与二次悬挂同构}
\label{g2:subsec:2403:07}
待撰写

\subsection{复周期性、球面系数及有限 CW 空间的计算}
\label{g2:subsec:2403:08}
待撰写
% 纲要条目排为编号小节标题，并进入目录与书签；说明段仍保留于纲要注释中。

% !TeX root = ../../Book2.tex
% 纲要标题：### 24.4 陈特征与指标问题
% 纲要位置：几何学计划纲要.md，第 1963 行。
% 正文状态：初稿；一般指标定理采用正文所列的准确外部输入。

% 纲要内容开始
% 
% * 有理陈特征和乘法
% * 陈特征的有理同构
% * 椭圆符号的紧支 K 类
% * 解析指标、拓扑指标与指标定理
% 
% ---
% 
% 纲要内容结束

\section{陈特征与指标问题}
\label{g2:sec:2404}

一张虚丛可以有许多不同的代表，而上同调类便于加减和计算。第21章的陈类已经把丛变成了上同调数据，不过总陈类把直和变成乘法，尚不能直接成为 \(K^0\) 到上同调环的环同态。本节先换一组适合加法与张量积的示性类，随后解释：为什么一个微分方程的最高阶部分也会确定一张虚丛，以及这份拓扑数据如何给出解空间之间的维数差。

\subsection{有理陈特征和乘法}
\label{g2:subsec:2404:01}
设 \(L\) 是复线丛，记 \(x=c_1(L)\)。线丛张量积满足 \(c_1(L\otimes L')=x+x'\)，所以指数式 \(e^x\) 正好把张量积变成乘法。若一张秩 \(r\) 的丛分裂为线丛之和，就把这些指数相加。分裂原理允许我们将这个想法用于一般复向量丛。

\begin{definition}\label{g2:def:chern-character}
	设 \(X\) 为有限 CW 空间，\(E\to X\) 为秩 \(r\) 的复向量丛。在\cref{g2:thm:splitting-principle}的分裂空间上记其陈根为 \(x_1,\ldots,x_r\)。定义\term[chen tezheng]{陈特征}[Chern character]
	\[
		\operatorname{ch}(E)\coloneqq\sum_{a=1}^r e^{x_a}
		=\sum_{j\ge0}\operatorname{ch}_j(E),\qquad
		\operatorname{ch}_0(E)=r,\quad
		\operatorname{ch}_j(E)=\frac1{j!}\sum_{a=1}^r x_a^j\in H^{2j}(X;\mathbb Q).
	\]
	这里各个对称多项式先写成 \(c_1(E),\ldots,c_r(E)\) 的多项式，再视为 \(X\) 上的类。若 \(X\) 不连通，零次项逐个连通分支取秩。
\end{definition}

这个定义中没有无穷级数的收敛问题：\(H^q(X;\mathbb Q)=0\) 当 \(q>\dim X\)，所以只有有限项可能非零。为何各项确实降到 \(X\)，也可以直接从 Newton 恒等式看出。令 \(p_j=\sum_a x_a^j\)、\(c_0=1\)，则
\[
	p_j-c_1p_{j-1}+c_2p_{j-2}-\cdots+(-1)^{j-1}c_{j-1}p_1+(-1)^j j c_j=0,
\]
其中 \(c_j=0\) 当 \(j>r\)。因此 \(p_j\) 是陈类的整数系数多项式，除以 \(j!\) 后给出唯一的有理类。前几项为
\begin{equation}\label{g2:eq:chern-character-low-terms}
	\operatorname{ch}(E)=r+c_1+\frac{c_1^2-2c_2}{2}
	+\frac{c_1^3-3c_1c_2+3c_3}{6}+\cdots.
\end{equation}
这些分母说明，有理系数从定义开始就有实际作用。例如在 \(\mathbb CP^2\) 上取 \(H=\gamma^*\)、\(h=c_1(H)\)，便有 \(\operatorname{ch}(H)=1+h+h^2/2\)。

\begin{proposition}\label{g2:prop:chern-character-ring}
	对有限 CW 空间 \(X\)，公式
	\[
		\operatorname{ch}([E]-[F])=\operatorname{ch}(E)-\operatorname{ch}(F)
	\]
	给出自然的含幺环同态
	\[
		\operatorname{ch}:K^0(X)\longrightarrow H^{\mathrm{even}}(X;\mathbb Q)
		\coloneqq\bigoplus_{j\ge0}H^{2j}(X;\mathbb Q).
	\]
	它与外积相容，且 \(\operatorname{ch}_j(E^*)=(-1)^j\operatorname{ch}_j(E)\)。
\end{proposition}
\begin{proof}
	先在同一个分裂空间上把 \(E,F\) 都分裂。\Cref{g2:thm:splitting-principle}保证其拉回在有理上同调中单射，所以在那里验证的等式能够降回 \(X\)。设根分别为 \(x_a,y_b\)。直和的根是两组根的并，故
	\[
		\operatorname{ch}(E\oplus F)=\sum_a e^{x_a}+\sum_b e^{y_b}.
	\]
	张量积则分裂为所有 \(L_a\otimes M_b\) 的直和，其根是 \(x_a+y_b\)，于是
	\[
		\operatorname{ch}(E\otimes F)
		=\sum_{a,b}e^{x_a+y_b}
		=\Bigl(\sum_a e^{x_a}\Bigr)\Bigl(\sum_b e^{y_b}\Bigr).
	\]
	这里根的次数为 \(2\)，杯积交换，有限截断的指数因此满足通常的乘法恒等式。平凡线丛的根为零，特征为 \(1\)。加法性使 \(\operatorname{ch}\) 经群完备化唯一延拓；若两对丛在稳定同构后相等，其陈特征之差也相等，因此虚丛的表示无关。上述乘法式随之延拓到虚丛。

	陈类与拉回相容，定义中的多项式也与拉回相容，故所得同态自然。对 \(X\times Y\) 上两个投影的拉回应用乘法式，就得到外积公式。最后，\(E^*\) 的根为 \(-x_a\)，其 \(j\) 次幂和是原来的 \((-1)^j\) 倍，给出对偶公式。
\end{proof}

对秩零类，零次分量消失。特别地，本章的 Bott 元满足
\begin{equation}\label{g2:eq:chern-bott-positive}
	\operatorname{ch}(\beta)=\operatorname{ch}([H]-1)=e^h-1=h
	\quad\text{于 }\widetilde H^*(S^2;\mathbb Q),
\end{equation}
因为 \(h^2=0\)。这既固定了接下来悬挂比较的正方向，也说明 Bott 元在上同调中对应二次生成元。

\subsection{陈特征的有理同构}
\label{g2:subsec:2404:02}
负一次 \(K\) 群由悬挂定义，因而其陈特征也应该由悬挂得到。记
\[
	s:H^q(X,A;\mathbb Q)\xrightarrow{\cong}
	\widetilde H^{q+1}(\Sigma(X/A);\mathbb Q)
\]
为与\secref{g2:sec:2402}锥坐标方向一致的上同调悬挂同构。将相对 \(K\) 类视为商空间上的约化类，可定义相对的偶陈特征。对 \(u\in K^{-1}(X,A)=\widetilde K^0(\Sigma(X/A))\)，定义
\[
	\operatorname{ch}^{-1}(u)\coloneqq s^{-1}\operatorname{ch}(u)
	\in H^{\mathrm{odd}}(X,A;\mathbb Q)
	\coloneqq\bigoplus_{j\ge0}H^{2j+1}(X,A;\mathbb Q).
\]
悬挂后的约化零次上同调为零，所以右端恰由正的奇次数组成。绝对情形采用 \(X/\varnothing=X_+\)，与本章的未约化 \(K^{-1}(X)\) 一致。

\begin{theorem}\label{g2:thm:rational-chern-character}
	对有限 CW 对 \((X,A)\)，陈特征给出自然同构
	\[
	\begin{aligned}
		K^0(X,A)\otimes_{\mathbb Z}\mathbb Q&\xrightarrow{\cong}H^{\mathrm{even}}(X,A;\mathbb Q),\\
		K^{-1}(X,A)\otimes_{\mathbb Z}\mathbb Q&\xrightarrow{\cong}H^{\mathrm{odd}}(X,A;\mathbb Q).
	\end{aligned}
	\]
	绝对情形的第一行是含幺环同构。
\end{theorem}
\begin{proof}
	证明按胞腔数归纳。需要先把球面上的比较与连接映射的相容性说明清楚，否则不能仅凭两边都具有长正合列就应用五引理。

	由\cref{g2:thm:complex-bott-periodicity,g2:prop:k-sphere-coefficients}，偶维球面 \(S^{2m}\) 的约化零次 \(K\) 群由 \(m\) 个 Bott 元的约化外积生成。式\eqref{g2:eq:chern-bott-positive}及外积公式把这个生成元送到 \(h_1\boxtimes\cdots\boxtimes h_m\)，即压缩积 \(S^2\wedge\cdots\wedge S^2\cong S^{2m}\) 的有向上同调生成元。另一种奇偶性的球面群为零。奇维球面悬挂一次成为偶维球面；按刚才的定义，负一次陈特征同样把相应生成元送到奇次上同调生成元。零维情形则由秩给出。故两种比较在每个球面及其有限楔和上都是有理同构。

	再看连接映射。\Cref{g2:thm:k-long-exact}中的连接由映射锥的商映射 \(X\cup_A CA\to\Sigma A\) 拉回给出；上同调的连接也是同一商映射的拉回，再经悬挂识别。陈特征与拉回相容，而我们用这同一个悬挂方向定义 \(\operatorname{ch}^{-1}\)，所以由负次数群组成的正合列与上同调列之间各个方块都交换。当用 Bott 周期性将次数折回偶、奇两类时，式\eqref{g2:eq:chern-bott-positive}说明乘 \(\beta\) 正对应两次上同调悬挂，因而这些方块仍然交换。这一步只使用已经证明的 Bott 周期性，并未用待证的有理同构来推断周期性。

	设 \(X^n\) 是 \(n\) 维骨架。相对商 \(X^n/X^{n-1}\) 是有限个 \(S^n\) 的楔和，故相对群上的比较已经是同构。对两个正合列取连续五项
	\[
		K^{j-1}(X^{n-1})\longrightarrow K^j(X^n,X^{n-1})
		\longrightarrow K^j(X^n)\longrightarrow K^j(X^{n-1})
		\longrightarrow K^{j+1}(X^n,X^{n-1}),
	\]
	其中次数按二周期解释。将此列张量 \(\mathbb Q\)，正合性保持：把每个元素写成分数后清除有限个分母，就把核像问题化回原来的整数正合列。其下方放置相同奇偶性的有理上同调列。由骨架归纳及球面计算，中间一项以外的四个竖箭头均为同构，五引理给出中间一项的同构。由有限零骨架出发，经过有限次便得到绝对结论。对任意有限 CW 对再取其长正合列，应用同样的五引理，得到相对结论。乘法相容性已由\cref{g2:prop:chern-character-ring}证明。
\end{proof}

因此，对有限 CW 空间，复 \(K\) 理论与普通上同调的差异包含两部分：有理化前的扭元，以及有理上同调中哪些类实际来自整数 \(K\) 类的整性条件。\Cref{g2:exm:k-moore-torsion}中的 Moore 空间给出第一种现象：其约化 \(K^0\) 群是 \(\mathbb Z/q\)，而正次数有理上同调全为零，陈特征看不到这些非零扭元。上面的 \(\mathbb CP^2\) 算例则说明，来自实际线丛的特征也可以带有分数系数。

\subsection{椭圆符号的紧支 K 类}
\label{g2:subsec:2404:03}
设 \(M\) 是闭光滑流形，\(E,F\to M\) 是有限秩复光滑向量丛，\(P:C^\infty(M,E)\to C^\infty(M,F)\) 是 \(m\ge1\) 阶线性微分算子。在局部坐标和丛标架下写成
\[
	P=\sum_{|\alpha|\le m}a_\alpha(x)\partial_x^\alpha.
\]
固定 Fourier 约定，以 \(\partial_{x_j}\) 对应 \(i\xi_j\)。其\term[zhufuhao]{主符号}[principal symbol]为
\[
	\sigma_m(P)(x,\xi)=\sum_{|\alpha|=m}a_\alpha(x)(i\xi)^\alpha:
	E_x\longrightarrow F_x.
\]
换坐标时，最高阶导数按余切变量作张量变换；导数落在换坐标或换标架的系数上时只贡献低阶项。因此这个公式给出 \(\mathrm T^*M\) 上由 \(\pi^*E\) 到 \(\pi^*F\) 的整体丛映射。对 \(t>0\) 有 \(\sigma_m(P)(x,t\xi)=t^m\sigma_m(P)(x,\xi)\)。

\begin{definition}\label{g2:def:elliptic-symbol-k-class}
	若对每个 \(x\in M\) 及非零 \(\xi\in\mathrm T_x^*M\)，\(\sigma_m(P)(x,\xi)\) 都是复线性同构，则称 \(P\) 为\term[tuoyuan suanzi]{椭圆算子}[elliptic operator]。取一个 Riemann 度量，以 \(D(\mathrm T^*M)\)、\(S(\mathrm T^*M)\) 表示单位闭余球丛及其边界。定义其符号类为
	\[
		[\sigma(P)]\coloneqq
		[\pi^*E,\pi^*F,\sigma_m(P)|_{S(\mathrm T^*M)}]
		\in K^0\bigl(D(\mathrm T^*M),S(\mathrm T^*M)\bigr).
	\]
\end{definition}

这里使用的正是\cref{g2:prop:k-triple-model}的三元组模型：两张丛在球丛内部不必同构，但在边界上由可逆主符号给出同构。将向量的长度径向压入单位球后，余切丛的一点紧化同胚于 \(D(\mathrm T^*M)/S(\mathrm T^*M)\)。因而这个群也写成
\[
	K_c^0(\mathrm T^*M)\coloneqq\widetilde K^0\bigl((\mathrm T^*M)^+\bigr),
\]
称为\term[jinzhi K lilun]{紧支 \(K\) 理论}[compactly supported K-theory]。下标 \(c\) 表达的是相对平凡化数据：在紧集以外，两个丛由指定同构相抵消。丛本身并不需要在零截面外消失。

\ProseParagraphBreak

还有一个有用的等价描述。在 \(\mathrm T^*M\) 上考虑两项复形
\[
	0\longrightarrow\pi^*E\xrightarrow{\ \sigma_m(P)\ }\pi^*F\longrightarrow0.
\]
它在零截面外正合，零截面又因 \(M\) 闭而紧，所以它代表一个紧支类。一般有限复形在紧集外正合，也可给出这样的类：取 Hermitian 度量，将偶次项与奇次项分别相加，以 \(\mathrm d+\mathrm d^*\) 从偶次和映到奇次和。正合处的拉普拉斯算子 \(\mathrm d\mathrm d^*+\mathrm d^*\mathrm d\) 可逆，所以这个奇偶映射可逆。两项复形时便恢复上述三元组。此模型及其乘积的详细构造可见\cite[第2节，第489--494页]{AtiyahSinger1968I}。

\begin{proposition}\label{g2:prop:symbol-class-invariance}
	闭光滑流形上复向量丛间椭圆微分算子的符号类不依赖定义单位余球时所取的度量。连续变形的可逆主符号给出相同的 \(K_c^0(\mathrm T^*M)\) 类；特别地，增加低于主阶数的项不改变该类。直和算子的符号类是原来两类之和。
\end{proposition}
\begin{proof}
	两个度量之间以正定二次型的凸组合连接。沿每条非零余切射线作径向重新缩放，便把各个单位余球丛连续识别；主符号的正齐次性说明重新缩放只乘一个正实数。这个因子可在正数中连接到 \(1\)，因此三元组类不变。主符号的同伦正是三元组允许的边界同构同伦。低阶项不进入 \(\sigma_m(P)\)，直和则给出主符号的块对角和，与三元组的加法一致。
\end{proof}

在圆周 \(M=\mathbb R/2\pi\mathbb Z\) 上，\(P=\mathrm d/\mathrm d\theta\) 作用于复值函数，主符号是 \(i\xi\)。单位余球有两个分支，边界同构分别是常数 \(i\) 与 \(-i\)。在这两个分支上各取一条从该常数到 \(1\) 的复非零道路，就将边界同构同伦为恒等；恒等同构延拓到整个单位余球丛，所以
\begin{equation}\label{g2:eq:circle-derivative-symbol-zero}
	[\sigma(\mathrm d/\mathrm d\theta)]=0\quad\text{于 }K_c^0(\mathrm T^*S^1).
\end{equation}
这是一次相对 \(K\) 类的计算，尚未数微分方程的解。

\subsection{解析指标、拓扑指标与指标定理}
\label{g2:subsec:2404:04}
现在从同一算子 \(P\) 出发，分别构造一个分析整数与一个拓扑整数。分析一侧必须先保证两个被相减的维数有限。所用的椭圆分析结论是：闭光滑 \(M\) 上的 \(m\) 阶椭圆算子，在 Sobolev 完备化后给出 Fredholm 算子
\[
	P:H^s(M,E)\longrightarrow H^{s-m}(M,F),
\]
其像闭，核与余核有限维，核及伴随核中的弱解均光滑，且指标与 \(s\) 无关。\(H^s\) 在整数 \(s\ge0\) 时可以理解为直到 \(s\) 阶弱导数平方可积的截面空间；完备化使“闭像”成为可用的分析性质。这里调用\cite[第6节，式(6.1)，第517--519页]{AtiyahSinger1968I}中的结论，Sobolev 截面与 Fredholm 分析分别安排在第三卷第11.2节、第13.2节，指标的几何应用见第三卷第17.4节的安排。

\begin{definition}\label{g2:def:analytic-index}
	设 \(M\) 闭且光滑，\(P:C^\infty(M,E)\to C^\infty(M,F)\) 为复向量丛间的椭圆微分算子。它的\term[jiexi zhibiao]{解析指标}[analytic index]为
	\[
		\operatorname{ind}_{\mathrm a}(P)\coloneqq
		\dim_{\mathbb C}\ker P-\dim_{\mathbb C}\operatorname{coker}P.
	\]
	这里余核可以在上述 Sobolev 空间中取商；由椭圆正则性，也可用光滑截面取商。选度量后，余核与形式伴随 \(P^*\) 的核自然对应。
\end{definition}

拓扑一侧要把余切丛上的紧支类送到一点上的 \(K^0(\mathrm{pt})=\mathbb Z\)。与普通上同调的 Gysin 映射相似，这需要法丛上的 Thom 类。复向量丛恰好具有一个可直接构造的 \(K\) 理论 Thom 类。

\begin{proposition}\label{g2:prop:complex-k-thom}
	设 \(V\to X\) 是有限 CW 空间上的秩 \(r\) 复向量丛。以 \(p:V\to X\) 为投影，\(v\in p^*V\) 为总空间上的典范截面。外代数复形
	\[
		0\longrightarrow p^*\Lambda^0V
		\xrightarrow{v\wedge}p^*\Lambda^1V
		\xrightarrow{v\wedge}\cdots
		\xrightarrow{v\wedge}p^*\Lambda^rV\longrightarrow0
	\]
	定义一个类 \(\lambda_V\in K_c^0(V)\)。乘这个类给出同构
	\[
		\Phi_V:K^j(X)\xrightarrow{\cong}K_c^j(V),\qquad
		a\longmapsto p^*a\cdot\lambda_V,\qquad j=0,-1.
	\]
	其相对版本也成立；若 \(X\) 是闭光滑流形上一张有限秩实光滑向量丛的总空间，\(V\to X\) 为复光滑向量丛，同样得到紧支版本 \(K_c^j(X)\cong K_c^j(V)\)。
\end{proposition}
\begin{proof}
	取 Hermitian 度量。记 \(\varepsilon(v)\) 为楔乘 \(v\)，\(\iota(v)=\varepsilon(v)^*\) 为其伴随；在每个纤维上
	\[
		\varepsilon(v)\iota(v)+\iota(v)\varepsilon(v)=|v|^2\operatorname{id}.
	\]
	所以 \(v\ne0\) 时，\(\iota(v)/|v|^2\) 是这个复形的收缩同伦，复形正合。等价地，\(\varepsilon(v)+\iota(v)\) 从偶次外幂和映到奇次外幂和，在零截面外可逆。由此得到紧支三元组。度量间的道路给出同伦，外代数构造与丛同构相容，故此类整体良定义。

	在一维复纤维上，复形是 \(\mathbb C\xrightarrow{z}\mathbb C\)。按\cref{g2:prop:k-triple-model}的粘合方向，原点盘的标架在无穷远盘中等于 \(z\) 倍标架。这是典范线丛 \(\gamma\to\mathbb CP^1\) 的过渡式；因此
	\begin{equation}\label{g2:eq:k-koszul-bott-sign}
		\lambda_{\mathbb C}=[\gamma]-1=[H]^{-1}-1=-\beta.
	\end{equation}
	最后一个等式使用 \(\beta^2=0\)。这也可以从本章粘接约定核对：正向 \(H\) 将无穷远盘的坐标乘 \(z\) 送入原点盘，而上述三元组将原点盘的坐标乘 \(z\) 送入无穷远盘，故二者互为逆线丛。

	对直和有带分次的外代数同构 \(\Lambda^*(V\oplus W)\cong\Lambda^*V\widehat\otimes\Lambda^*W\)，微分对应张量总微分。因此 \(\lambda_{V\oplus W}=\lambda_V\boxtimes\lambda_W\)（在共同底空间上取相应拉回）。在平凡秩 \(r\) 丛上，乘 \(\lambda_V\) 就是乘 \((-1)^r\beta^{\boxtimes r}\)，由复 Bott 周期性是同构。

	一般 \(X\) 按有限骨架过滤。将 \(V\) 限制到一个胞腔的特征圆盘上，它平凡；对于一对相邻骨架，相对 Thom 空间的商是这些圆盘上相对 Thom 空间的有限楔和，因而相对映射正是刚才证明的平凡丛版本。Thom 类由整体外代数复形给出，限制到骨架和边界时保持同一复形，故各个长正合列方块交换。从零骨架开始应用五引理，得到结论。对 \((X,A)\) 用相对骨架作同样的归纳，证明相对版本。

	最后，若底空间 \(X\) 是闭光滑流形上一张有限秩实光滑向量丛的总空间，以半径递增的闭球丛穷尽它。每个紧集都包含在某个球丛的内部；紧支 \(K\) 类因而来自相应球丛及边界的相对群。球丛与边界具有有限 CW 对的同伦型，上一段的相对同构可用，并与扩大半径相容。取这些群的直极限，便得到所述紧支版本。此时外代数复形本身只在纤维方向具有紧支；与底空间的紧支类相乘后，所得复形在总空间中才具有紧支。
\end{proof}

这里采用的是原始指标论文的外代数 \(V\) 约定，见\cite[第2节，第493--494页]{AtiyahSinger1968I}。式\eqref{g2:eq:k-koszul-bott-sign}说明它与本章正陈根的 Bott 元相差一个负号；下述推前和末端整数识别始终使用同一套 \(\lambda\) 约定。

\ProseParagraphBreak

取光滑嵌入 \(i:M\hookrightarrow\mathbb R^N\)，记其实法丛为 \(\nu\)。度量把 \(\mathrm T^*M\) 识别为 \(\mathrm TM\)，而切映射给出嵌入 \(\mathrm Ti:\mathrm TM\hookrightarrow\mathrm T\mathbb R^N\)。其法丛同构于从 \(M\) 拉回的 \(\nu\oplus\nu\)：第一份来自底空间的法位移，第二份来自切向量的法分量。这个分解可由一个联络写出，改变联络只产生同伦。按
\[
	(a,b)\longmapsto a+ib
\]
给两份法方向赋予复结构，就得到 \(p^*(\nu\otimes_{\mathbb R}\mathbb C)\)。它的复秩为 \(N-\dim M\)。因此此处的复 Thom 类来自两份相同实法丛的组合。

\ProseParagraphBreak

先沿这个复法丛作 Thom 映射，再将管状邻域中的紧支类延为环境中零，得到
\[
	(\mathrm Ti)_!:K_c^0(\mathrm TM)\longrightarrow
	K_c^0(\mathrm T\mathbb R^N).
\]
“延为零”由一点紧化之间的压缩映射实现：环境中管状邻域以外的部分送到无穷远点，故其方向是 \(K\) 类的拉回。将 \((x,v)\in\mathrm T\mathbb R^N\) 写成 \(x+iv\in\mathbb C^N\)，前一命题及 Bott 周期性给出
\[
	K_c^0(\mathrm T\mathbb R^N)=K_c^0(\mathbb C^N)
	=\mathbb Z\,\lambda_{\mathbb C^N},\qquad
	\lambda_{\mathbb C^N}=\lambda_{\mathbb C}^{\boxtimes N}=(-1)^N\beta^{\boxtimes N}.
\]
令 \(b_N\) 把 \(\lambda_{\mathbb C^N}\) 送到 \(1\)。于是\term[tuopu zhibiao]{拓扑指标}[topological index]是复合
\begin{equation}\label{g2:eq:topological-index-pushforward}
	\operatorname{ind}_{\mathrm t}:K_c^0(\mathrm T^*M)
	\xrightarrow{\ \cong\ }K_c^0(\mathrm TM)
	\xrightarrow{\ (\mathrm Ti)_!\ }K_c^0(\mathbb C^N)
	\xrightarrow{\ b_N\ }\mathbb Z.
\end{equation}
这里既固定了法丛中“底法方向在前、纤维法方向在后”的次序，也固定了末端生成元，因而不会因混用 Bott 元的符号改变指标。

\ProseParagraphBreak

这个整数与嵌入和辅助度量无关。度量的变化给出一族识别，管状邻域的变化给出同伦的压缩映射；Thom 类在这些同伦下自然。增加一个平凡环境实方向时，推前额外乘 \(\lambda_{\mathbb C}\)，而末端 \(b_{N+1}\) 正好消去这一因子，故整数不变。对两个嵌入 \(i_0:M\to\mathbb R^{N_0}\)、\(i_1:M\to\mathbb R^{N_1}\)，在乘积环境中用
\[
	x\longmapsto\bigl(\cos t\,i_0(x),\sin t\,i_1(x)\bigr),
	\qquad 0\le t\le\pi/2,
\]
连接其稳定化。每个中间映射至少有一个非零系数的嵌入分量，因而仍是嵌入；沿这族嵌入的法丛及管状邻域给出同伦的推前。末端交换环境实坐标只在复化后作复线性变换，保持 \(\lambda\) 类。上述构造与独立性的完整拓扑处理见\cite[第3节，第497--499页]{AtiyahSinger1968I}。

\begin{theorem}[Atiyah--Singer 指标定理]\label{g2:thm:atiyah-singer-topological-index}
	设 \(M\) 为闭光滑流形，\(E,F\to M\) 为有限秩复光滑向量丛，\(P:C^\infty(M,E)\to C^\infty(M,F)\) 为椭圆微分算子。则
	\[
		\operatorname{ind}_{\mathrm a}(P)
		=\operatorname{ind}_{\mathrm t}([\sigma(P)]),
	\]
	其中右端使用式\eqref{g2:eq:topological-index-pushforward}的复法丛推前及 \(\lambda_{\mathbb C^N}\mapsto1\) 的归一化。
\end{theorem}

这个结论称为\term[Atiyah-Singer-zhibiao-dingli]{Atiyah--Singer 指标定理}[Atiyah--Singer index theorem]。\footnote{Michael Francis Atiyah（1929--2019），英国数学家，参与建立拓扑 \(K\) 理论，并研究几何与数学物理之间的联系。Isadore Manuel Singer（1924--2021），美国数学家，研究微分几何与泛函分析。}这里引用\cite[第6节，定理(6.7)，第520页]{AtiyahSinger1968I}的结论，取其等变陈述中的作用群为平凡群。完整证明还要建立椭圆算子的解析指标在这些拓扑推前下的相容性；原文第8、9节完成这一环节。解析指标从有限维核与余核得到整数，拓扑指标从符号的紧支类及法丛得到整数，定理把这两种构造联系起来。

\begin{example}\label{g2:exm:circle-derivative-index}
	对圆周上的 \(P=\mathrm d/\mathrm d\theta\)，周期解 \(Pf=0\) 恰是常函数，核为一维。方程 \(Pf=g\) 有周期光滑解，当且仅当
	\[
		\int_0^{2\pi}g(\theta)\,\mathrm d\theta=0.
	\]
	必要性来自对导数积分；充分性由 \(f(\theta)=\int_0^\theta g(t)\,\mathrm dt\) 给出，零积分保证周期性，且 \(g\) 的周期性保证各阶导数在接缝吻合。因此余核由平均值给出，也是一维，解析指标为 \(1-1=0\)。另一方面，式\eqref{g2:eq:circle-derivative-symbol-zero}已算出符号类为零，拓扑指标随之为零。这个例子中，核与余核都非零，二者的差却只由符号的拓扑类决定。
\end{example}
% 纲要条目排为编号小节标题，并进入目录与书签；说明段仍保留于纲要注释中。


\begin{chaptersummary}
	K 群以虚丛保留直和的加法关系，以张量积保留乘法。紧空间上的补丛把虚丛写成 \([E]-[\underline{\mathbb C}^N]\)，相对类则同时记录子空间上的一个稳定同构。映射锥使这份边界数据成为一个约化类，悬挂与余纤维列由此给出长正合列和切除。

	\ProseParagraphBreak

	复 Bott 周期性的证明把连续粘接函数依次化为 Laurent 多项式和一次矩阵。谱投影将一次矩阵分成圆内与圆外两类子丛：圆内部分贡献 \(H\)，圆外部分贡献常丛。对多项式次数、负幂及同伦选择的比较保证逆映射良定义，从而得到 \(K^0(X\times S^2)=K^0(X)\oplus K^0(X)\beta\) 以及 \(\beta^2=0\)。取约化部分后，二次悬挂保持复 K 群；球面系数、六项正合列和逐胞腔计算随之得到。

	\ProseParagraphBreak

	陈特征把直和与张量积分别变成上同调的加法与乘法，并在有限 CW 空间上给出有理同构。它揭示有理信息，也解释扭元为何可能不被示性类看见。椭圆算子的可逆主符号给出余切丛上的紧支 K 类；复 Thom 类和嵌入把它推到一个整数。解析指标由核与余核定义，拓扑指标由这个推前定义，Atiyah--Singer 指标定理断言二者相等。
\end{chaptersummary}

\BookExercises

\begin{exercise}\label{g2:ex:k-disconnected-reduction}
	设 \(X\) 是 \(m\ge2\) 个点组成的离散空间，取第一个点为基点。求 \(K^0(X)\) 的环结构及 \(\widetilde K^0(X)\)。再求“在每一个点虚秩都为零”的类组成的子群，比较它与约化群。
\end{exercise}
\begin{solution}
	一张丛由每个点上的复向量空间决定，同构类由 \((r_1,\ldots,r_m)\in\mathbb N^m\) 唯一确定。直和逐坐标相加，所以群完备化为 \(\mathbb Z^m\)；张量积逐纤维进行，乘法也逐坐标进行：
	\[
		(a_1,\ldots,a_m)(b_1,\ldots,b_m)=(a_1b_1,\ldots,a_mb_m).
	\]
	单位元为 \((1,\ldots,1)\)。基点取秩是第一坐标投影，故 \(\widetilde K^0(X)=\{0\}\times\mathbb Z^{m-1}\)。在每点虚秩均为零则要求全部坐标都为零，所得子群为零。约化群只消去所选基点处的秩；例如 \((0,1,0,\ldots,0)\) 是约化群中的非零元。
\end{solution}

\begin{exercise}\label{g2:ex:hopf-dual-sum-trivialization}
	设 \(H\to S^2\) 是本章的 Hopf 线丛。用粘合矩阵给出 \(H\oplus H^*\cong\varepsilon^2\) 的具体同伦证明，并计算 \([H^*]-1\) 与 \([H]-1\) 的关系。解释为什么这个关系不意味着 \([H^*]=-[H]\)。
\end{exercise}
\begin{solution}
	按本章赤道约定，直和的粘合矩阵为 \(\operatorname{diag}(z,z^{-1})\)，\(|z|=1\)。下式在整个参数区间取值于 \(\operatorname{SU}(2)\)：
	\[
		G_t(z)=\begin{pmatrix}\cos t\,z&-\sin t\\\sin t&\cos t\,\overline z\end{pmatrix},
		\qquad0\le t\le\pi/2.
	\]
	两列正交且为单位向量，行列式为 \(\cos^2t+\sin^2t=1\)。起点是所给粘合函数，终点为常矩阵 \(J=\left(\begin{smallmatrix}0&-1\\1&0\end{smallmatrix}\right)\)。再用实旋转矩阵把 \(J\) 连到 \(I\)，便得到粘合函数的零伦。\cref{g2:thm:clutching-classification}于是给出实际秩二丛的平凡化类型。

	在 \(K^0(S^2)\) 中，这说明 \([H]+[H^*]=2\)，所以 \([H^*]-1=-([H]-1)\)。反之，\([H^*]\) 的虚秩为 \(1\)，\(-[H]\) 的虚秩为 \(-1\)，二者不相等。负元关系成立的是已经减去秩的两项。
\end{solution}

\begin{exercise}\label{g2:ex:k-pullback-two-sphere-degree}
	设 \(f:S^2\to S^2\) 的映射度为整数 \(d\)，记 \(\eta=[H]-1\)。求 \(f^*:K^0(S^2)\to K^0(S^2)\) 在加法基 \(1,\eta\) 下的矩阵，并分别求其核与余核。最后说明 \(f^*H\) 与哪张线丛同构。
\end{exercise}
\begin{solution}
	拉回保持平凡线丛，所以 \(f^*1=1\)。由第一陈类的自然性，\(c_1(f^*H)=f^*h=dh\)，故式\eqref{g2:eq:k-two-sphere-additive}给出 \(f^*\eta=d\eta\)。因此所求矩阵为
	\[
		\begin{pmatrix}1&0\\0&d\end{pmatrix}.
	\]
	当 \(d\ne0\) 时，核为零，余核为 \(\mathbb Z/|d|\mathbb Z\)；当 \(d=0\) 时，核为 \(\mathbb Z\eta\)，余核为 \(\mathbb Z\)。秩一复丛由第一陈类分类，所以 \(f^*H\cong H^{\otimes d}\)，负张量幂按对偶理解。这个同构由实际线丛分类得到，比只写出它在 \(K^0\) 中的等式更精确。
\end{solution}

\begin{exercise}\label{g2:ex:k-external-product-visible-chern}
	在 \(S^2\times S^2\) 上，记 \(L=p_1^*H\)、\(M=p_2^*H\)。证明 \(([L]-1)([M]-1)\ne0\)，但该类在 \(S^2\times\{*\}\) 和 \(\{*\}\times S^2\) 上的限制均为零。可直接使用第21章的陈类和公式，无需使用 K 理论乘积定理。
\end{exercise}
\begin{solution}
	令 \(a=p_1^*h\)、\(b=p_2^*h\)。普通整数上同调的 K\"unneth 公式给出 \(a^2=b^2=0\)，而 \(ab\) 在乘积基本类上的求值为 \(1\)，故非零。展开虚丛乘积为
	\[
		([L]-1)([M]-1)=[(L\otimes M)\oplus\varepsilon^1]-[L\oplus M].
	\]
	若它为零，两张秩二丛应当稳定同构，故具有相同陈类。可是
	\[
		c((L\otimes M)\oplus\varepsilon^1)=1+a+b,
		\qquad c(L\oplus M)=(1+a)(1+b)=1+a+b+ab.
	\]
	二次陈类之差为非零的 \(ab\)，矛盾。在第一张因子球面上，\(M\) 限制为平凡线丛；在第二张因子球面上，\(L\) 限制为平凡线丛。因此乘积在两者上均为零。这个类记录两因子共同参与的信息，分别限制到各因子会失去它。
\end{solution}

\begin{exercise}\label{g2:ex:k-retraction-torus}
	设有限 CW 对 \((X,A)\) 的包含有连续收缩映射 \(r:X\to A\)，即 \(r|_A=\operatorname{id}_A\)。证明对每个 \(q\ge0\)，
	\[
		K^{-q}(X)\cong K^{-q}(X,A)\oplus K^{-q}(A).
	\]
	将此用于 \(X=T^2=S^1\times S^1\)、\(A=S^1\times\{1\}\)，计算两者的零次、负一次相对群及 \(K^0(T^2),K^{-1}(T^2)\)。
\end{exercise}
\begin{solution}
	由 \(i^*r^*=\operatorname{id}\)，限制 \(i^*\) 在每个次数上满射。长正合列中进入相对群的连接映射因前一次限制满射而为零，所以
	\[
		0\longrightarrow K^{-q}(X,A)\xrightarrow{j}K^{-q}(X)
		\xrightarrow{i^*}K^{-q}(A)\longrightarrow0
	\]
	短正合，并以 \(r^*\) 为截面。明确的分解把 \(u\) 写成 \((u-r^*i^*u)+r^*i^*u\)，其中第一项在 \(j\) 的像内且有唯一相对原像。

	对环面，取向第一因子的投影作 \(r\)。标准胞腔分解有两条边 \(a,b\) 和一张沿 \(aba^{-1}b^{-1}\) 附着的二胞腔。压掉 \(A\) 就压掉 \(a\)，附着字变为 \(bb^{-1}\)，它零伦。由\cref{g2:prop:homotopic-attaching-maps}，
	\[
		T^2/A\simeq S^1\vee S^2.
	\]
	本章的球面计算给出 \(\widetilde K^0(S^1)=0\)、\(\widetilde K^0(S^2)=\mathbb Z\)，以及 \(\widetilde K^{-1}(S^1)=\mathbb Z\)、\(\widetilde K^{-1}(S^2)=0\)。因此
	\[
		K^0(T^2,A)\cong\mathbb Z,\qquad K^{-1}(T^2,A)\cong\mathbb Z.
	\]
	又 \(K^0(A)=\mathbb Z\)、\(K^{-1}(A)=\mathbb Z\)，代入刚证明的分裂，得到 \(K^0(T^2)\cong\mathbb Z^2\)、\(K^{-1}(T^2)\cong\mathbb Z^2\)。这是加性群的计算，乘法还须保留两个方向及其分次符号。
\end{solution}

\begin{exercise}\label{g2:ex:k-mapping-torus-degree}
	设 \(f:S^2\to S^2\) 为保持基点的度数 \(d\in\mathbb Z\) 的连续映射，定义映射环面
	\[
		T_f=(S^2\times[0,1])/\bigl((x,1)\sim(f(x),0)\bigr).
	\]
	用两段圆柱及相对正合列计算 \(K^0(T_f)\)、\(K^{-1}(T_f)\)，特别求出 \(d=-1\) 时的结果。说明计算是否要求 \(T_f\to S^1\) 是纤维丛。
\end{exercise}
\begin{solution}
	给 \(S^2\) 取一张零胞腔和一张二胞腔，保持基点的 \(f\) 就是胞腔映射，故 \(T_f\) 有有限 CW 结构。记 \(U\) 为参数区间 \([0,1/3]\cup[2/3,1]\) 的像，\(V\) 为 \([1/4,3/4]\) 的像。细分圆柱后，这些集合形成子复形覆盖。\(U\) 由接缝处的映射柱及一段领圈组成，形变收缩到 \(t=0\) 的目标球面；\(V\) 也形变收缩到一张球面。交集有两条圆柱，故
	\[
		U\simeq S^2,\qquad V\simeq S^2,\qquad U\cap V\simeq S^2\sqcup S^2.
	\]
	按靠近 \(t=0\)、靠近 \(t=1\) 的次序排列交集分支，进入 \(U\) 的两张映射在这些识别下分别为恒等与 \(f\)，进入 \(V\) 的两张映射均为恒等。

	这里使用的 Mayer--Vietoris 序列来自\cref{g2:thm:k-long-exact}：对 \((T_f,U)\) 与 \((V,U\cap V)\) 写长正合列，利用切除识别相对项；两条限制映射作差，逐项正合性便给出
	\[
		\cdots\longrightarrow K^i(T_f)\longrightarrow K^i(U)\oplus K^i(V)
		\xrightarrow{D_i}K^i(U\cap V)\longrightarrow K^{i+1}(T_f)\longrightarrow\cdots.
	\]
	次数按本章的二周期约定理解，\(K^1=K^{-1}\)。例如在中间项，差为零表示两份限制相同；先用一对的正合性把它们的差提升到相对项，再通过切除移到另一对，即得到整体原像。对连接项同样使用相对列的核像等式，得到相邻两处的正合性。这正是两段粘合所需的序列。

	在上述三个同伦模型中，差映射为
	\[
		D_i(a,b)=(a-b,f^*a-b).
	\]
	把第二个输出减去第一个输出，并以 \((a,a-b)\) 作为输入坐标，就将这张映射化为一份恒等映射与 \(f^*-1\) 的直和。因此其核、余核分别等于 \(1-f^*\) 的核、余核。由于 \(K^{-1}(S^2)=0\)，正合列给出
	\[
		K^0(T_f)\cong\ker(1-f^*:K^0(S^2)\to K^0(S^2)),\qquad
		K^{-1}(T_f)\cong\operatorname{coker}(1-f^*).
	\]
	以 \(1,\beta\) 为 \(K^0(S^2)=\mathbb Z\oplus\mathbb Z\beta\) 的基，\(f^*1=1\)。又 \(c_1:\widetilde K^0(S^2)\to H^2(S^2;\mathbb Z)\) 为同构，而度数 \(d\) 的映射把二次上同调乘以 \(d\)，故 \(f^*\beta=d\beta\)。所以
	\[
	\begin{aligned}
		K^0(T_f)&\cong\mathbb Z\oplus\ker(1-d:\mathbb Z\to\mathbb Z),\\
		K^{-1}(T_f)&\cong\mathbb Z\oplus\operatorname{coker}(1-d:\mathbb Z\to\mathbb Z).
	\end{aligned}
	\]
	当 \(d=-1\) 时，结果为 \(K^0(T_f)=\mathbb Z\)、\(K^{-1}(T_f)\cong\mathbb Z\oplus\mathbb Z/2\)。当 \(d=1\) 时，两群均为 \(\mathbb Z^2\)；其他次数时，前者为 \(\mathbb Z\)，后者为 \(\mathbb Z\oplus\mathbb Z/|1-d|\)。整个计算只用了映射柱、同伦型与粘合正合性，未要求 \(f\) 为同胚，也未要求 \(T_f\to S^1\) 是纤维丛。
\end{solution}

\begin{exercise}\label{g2:ex:k-suspension-square-zero}
	设 \(Y\) 是带零胞腔基点的有限 CW 空间。证明 \(K^0(\Sigma Y)\) 中的约化理想满足
	\[
		\widetilde K^0(\Sigma Y)\cdot\widetilde K^0(\Sigma Y)=0.
	\]
	解释这为什么并不要求 \(\widetilde K^0(\Sigma Y)\) 本身为零。
\end{exercise}
\begin{solution}
	乘法是外积沿对角映射的拉回。对约化类，它只取决于约化对角
	\[
		\overline\Delta:\Sigma Y\longrightarrow\Sigma Y\wedge\Sigma Y.
	\]
	将悬挂点写为 \([t,y]\)，\(0\le t\le1\)。在乘积 \(\Sigma Y\times\Sigma Y\) 中，令
	\[
		a(t)=\min(2t,1),\qquad b(t)=\max(2t-1,0).
	\]
	映射 \([t,y]\mapsto([a(t),y],[b(t),y])\) 落在楔和中：前半段第二坐标为基点，后半段第一坐标为基点。把两个参数函数分别沿直线从 \(t\) 变到 \(a(t),b(t)\)，就得到从普通对角到此映射的同伦；端点及 \(y\) 为基点时始终取基点，所以同伦在悬挂商上良定义。再投影到压缩积，便将 \(\overline\Delta\) 零伦化。

	因此对任意约化类 \(u,v\)，有 \(uv=\overline\Delta^*(u\boxtimes v)=0\)。但单个约化类可以非零，例如取 \(Y=S^1\)，则 \(\Sigma Y\cong S^2\)，其 Bott 元生成 \(\widetilde K^0(S^2)\cong\mathbb Z\)。这里消失的是两个约化类的乘积，而不是加性群。
\end{solution}

\begin{exercise}\label{g2:ex:k-family-spectral-bundle}
	在第一份 \(S^2=\mathbb CP^1\) 上，把 \(H=\gamma^*\) 实现为平凡丛 \(\underline{\mathbb C}^2\) 的一个线子丛，记相应正交投影为 \(P_x\)。令
	\[
		C_x=2(I-P_x),\qquad f_x(z)=zI+C_x\quad(|z|=1),
	\]
	并由 \(f\) 在 \(S^2\times S^2\) 上粘接一个秩二丛 \(V\)。若 \(\beta_1,\beta_2\) 是两因子上 Bott 元的拉回，求 \([V]\)。证明：虽然每个固定 \(x\) 上的限制 \(V|_{\{x\}\times S^2}\) 都同构于 \(\underline{\mathbb C}\oplus H\)，\(V\) 仍不与 \(W=\underline{\mathbb C}\oplus\operatorname{pr}_2^*H\) 稳定同构。
\end{exercise}
\begin{solution}
	取 \(\mathbb C^2\) 的标准 Hermitian 内积。\(\gamma^\perp\) 的第一陈类为 \(-c_1(\gamma)=c_1(H)\)，由球面上线丛的分类，它同构于 \(H\)，故所要求的线子丛确实存在。\(C\) 在 \(H\) 上为零，在其正交补 \(H^\perp\) 上为乘以2；\(zI+C\) 的两个特征值分别为 \(z\) 与 \(z+2\)，在单位圆上均不为零。

	在 \(H\) 上，粘接函数 \(z\) 给出 \(H\boxtimes H\)；在 \(H^\perp\) 上，\(2+sz\)（\(0\le s\le1\)）始终非零，把粘接函数变为常数2，再改换一个半球上的平凡化即可去掉这一常数。于是
	\[
		V\cong\operatorname{pr}_1^*(H^\perp)\oplus(H\boxtimes H).
	\]
	利用 \([H^\perp]=2-[H]=1-\beta_1\)，得到
	\[
		[V]=(1-\beta_1)+(1+\beta_1)(1+\beta_2)
		=2+\beta_2+\beta_1\beta_2.
	\]
	固定 \(x\) 时，两条来自第一因子的线都成为固定的一维向量空间，所以限制确为 \(\underline{\mathbb C}\oplus H\)。但是 \([W]=2+\beta_2\)，两者之差为 \(\beta_1\beta_2\)。由\cref{g2:thm:k-basic-product}，\(K^0(S^2\times S^2)=K^0(S^2)\oplus K^0(S^2)\beta_2\)，而 \(\beta_1\ne0\)，故这个差非零。稳定同构会给出相同的 K 类，因此不可能成立。这个例子说明，逐点计算绕数还没有记录单位圆内谱空间怎样随参数组成子丛。
\end{solution}

\begin{exercise}\label{g2:ex:k-polynomial-winding}
	设 \(q(z)=z^2-3z+1\)，\(m\) 是非负整数。计算单位圆上粘接函数 \(z^{-m}q(z)\) 所定义的线丛的 K 类。写出正文给出的 \(q\) 的一次矩阵线性化，并求其单位圆内谱子丛的秩。说明多项式次数为何不足以决定该秩。
\end{exercise}
\begin{solution}
	两个根是 \(r_-=(3-\sqrt5)/2\) 与 \(r_+=(3+\sqrt5)/2\)，满足 \(0<r_-<1<r_+\)。在单位圆上，\(z-sr_-\) 始终不为零，把 \(z-r_-\) 同伦到 \(z\)；\(sz-r_+\) 也始终不为零，把 \(z-r_+\) 同伦到常数 \(-r_+\)。因此 \(q(z)\) 的粘接同伦于 \(-r_+z\)。去掉非零常数后，\(z^{-m}q(z)\) 的线丛就是 \(H^{1-m}\)，所以
	\[
		[C(\mathbb C,z^{-m}q)]=1+(1-m)\beta.
	\]
	这里用了 \(\beta^2=0\)，因而整数幂均满足 \([H]^j=1+j\beta\)。

	按\cref{g2:eq:k-linearization-matrix}，可以取
	\[
		L_2(q)(z)=
		\begin{pmatrix}
			1&-z&0\\
			0&1&-z\\
			1&-3&1
		\end{pmatrix}.
	\]
	线性化公式给出 \([C(\mathbb C^3,L_2(q))]=2+[H]=3+\beta\)。另一方面，谱分解给出 \(3+\operatorname{rank}S(L_2(q))\,\beta\)。由 \(\widetilde K^0(S^2)=\mathbb Z\beta\)，比较系数即得 \(\operatorname{rank}S(L_2(q))=1\)。\(q\) 虽是二次多项式，却只有一个根位于单位圆内；圆外因子可同伦为非零常数，不贡献 Bott 元。比如 \(z^2\) 同样是二次多项式，所给类却为 \(1+2\beta\)，其对应谱秩为2。
\end{solution}

\begin{exercise}\label{g2:ex:chern-endomorphism-rank-two}
	设 \(E\to X\) 是有限 CW 空间上的秩二复向量丛，\(c_1(E)=0\)，且 \(\dim X\le8\)。计算 \(\operatorname{ch}(\operatorname{End}E)\)，并说明其中为何没有次数为 \(2\) 和 \(6\) 的分量。
\end{exercise}
\begin{solution}
	在分裂空间上设 \(E\) 的根为 \(x,y\)。由 \(c_1(E)=0\) 得 \(y=-x\)，从而 \(c_2(E)=xy=-x^2\)。因为 \(\operatorname{End}E=E^*\otimes E\)，它的四个根是 \(0,0,2x,-2x\)。于是
	\[
		\operatorname{ch}(\operatorname{End}E)=2+e^{2x}+e^{-2x}
		=4+4x^2+\frac43x^4
		=4-4c_2(E)+\frac43c_2(E)^2.
	\]
	超过次数 \(8\) 的项因空间维数而为零。次数 \(2\) 和 \(6\) 对应根的一次和三次幂和；成对的 \(2x,-2x\) 在奇数次幂中抵消。也可用 \(\operatorname{End}E\cong(\operatorname{End}E)^*\) 及陈特征的对偶公式，看出所有 \(\operatorname{ch}_{2j+1}\) 在有理系数下消失。分裂空间的上同调拉回单射，故等式在 \(X\) 上成立。
\end{solution}

\begin{exercise}\label{g2:ex:chern-lens-torsion-line}
	设 \(q\ge2\)，\(X=M_q=S^1\cup_{\deg q}D^2\)。取复线丛 \(L\to X\)，使 \(c_1(L)\) 为 \(H^2(X;\mathbb Z)\cong\mathbb Z/q\) 的一个生成元。
	\begin{enumerate}
		\item 证明 \(L^{\otimes q}\) 平凡，但 \(L\) 不稳定平凡。
		\item 证明 \(u=[L]-1\in\widetilde K^0(X)\) 的加法阶恰为 \(q\)，而 \(\operatorname{ch}(u)=0\)。区分第一问的张量幂与这里的加法倍数。
	\end{enumerate}
\end{exercise}
\begin{solution}
	1. \Cref{g2:prop:line-bundle-cohomology-classification}给出复线丛的上同调分类，第21章将其二次分类类识别为第一陈类；线丛张量积的第一陈类相加。因此 \(c_1(L^{\otimes q})=q c_1(L)=0\)，故 \(L^{\otimes q}\) 平凡。若 \(L\oplus\varepsilon^r\cong\varepsilon^{r+1}\)，取行列式线丛便得到 \(L\cong\varepsilon^1\)，这与 \(c_1(L)\ne0\) 矛盾。

	2. 若 \(nu=0\)，稳定等式意味着 \(L^{\oplus n}\) 与秩 \(n\) 平凡丛稳定同构（先取 \(n>0\)）。对这个直和取第一陈类，得到 \(n c_1(L)=0\)，所以 \(q\mid n\)。另一方面，\cref{g2:exm:k-moore-torsion}已经计算 \(\widetilde K^0(M_q)\cong\mathbb Z/q\)，所以 \(qu=0\)。二者合起来说明 \(u\) 的阶正好是 \(q\)。

	\(c_1(L)\) 是扭元，进入有理上同调后为零。因此 \(\operatorname{ch}(L)=e^{c_1(L)}=1\)，\(\operatorname{ch}(u)=0\)。第一问给出乘法关系 \([L]^q=1\)，第二问给出加法关系 \(q([L]-1)=0\)；后一个关系在本题中来自 \(K\) 群的计算，不能只把前一个等式中的幂换成倍数。
\end{solution}

\begin{exercise}\label{g2:ex:circle-derivative-zeroth-order}
	在复值周期函数上考虑 \(P_a=\mathrm d/\mathrm d\theta+a\)，其中 \(a\in\mathbb C\) 为常数，圆周周期为 \(2\pi\)。求其核与余核的维数，证明指标恒为零，并解释当 \(a\) 经过 \(i\mathbb Z\) 时，解空间维数的跳跃为何不违背指标的稳定性。
\end{exercise}
\begin{solution}
	齐次方程的解为 \(f(\theta)=Ce^{-a\theta}\)。它具有周期 \(2\pi\) 当且仅当 \(e^{-2\pi a}=1\)，即 \(a\in i\mathbb Z\)。所以核在这些参数处为一维，在其他参数处为零。

	对非齐次方程，积分因子给出
	\[
		f(\theta)=e^{-a\theta}\left(C+\int_0^\theta e^{at}g(t)\,\mathrm dt\right).
	\]
	周期条件等价于
	\[
		(e^{2\pi a}-1)C=\int_0^{2\pi}e^{at}g(t)\,\mathrm dt.
	\]
	若 \(a\notin i\mathbb Z\)，对任意 \(g\) 都唯一确定 \(C\)，故余核为零。若 \(a\in i\mathbb Z\)，可解条件是右端为零；这是一条非零的线性条件，例如取 \(g(t)=e^{-at}\) 时积分为 \(2\pi\)。所以余核为一维。由微分方程反复求导可知，在函数值匹配后各阶导数也周期匹配，所得解确为光滑周期函数。

	因此两种情况下核、余核的维数分别为 \((0,0)\) 与 \((1,1)\)，指标始终为零。参数经过 \(i\mathbb Z\) 时两者同时跳跃，差保持不变。所有 \(P_a\) 的主符号都是 \(i\xi\)，由式\eqref{g2:eq:circle-derivative-symbol-zero}，符号类恒为零，也给出同一拓扑指标。
\end{solution}

\begin{exercise}\label{g2:ex:adjoint-symbol-inverse}
	设 \(M\) 为闭光滑流形，\(P:C^\infty(M,E)\to C^\infty(M,F)\) 为复向量丛间的椭圆微分算子。选 Riemann 度量及 Hermitian 丛度量，以积分恒等式 \(\langle Pu,v\rangle_{L^2}=\langle u,P^*v\rangle_{L^2}\) 定义形式伴随。证明
	\[
		[\sigma(P^*)]=-[\sigma(P)]\quad\text{于 }K_c^0(\mathrm T^*M),
		\qquad\operatorname{ind}_{\mathrm a}(P^*)=-\operatorname{ind}_{\mathrm a}(P).
	\]
	分别从三元组和核、余核说明这两个负号的来源。
\end{exercise}
\begin{solution}
	分部积分把每个最高阶偏导数变为其负号，系数同时取 Hermitian 伴随；导数落在系数或体积密度上产生的都是低阶项。在采用 \(i\xi\) 的主符号约定下，这正给出 \(\sigma_m(P^*)=\sigma_m(P)^*\)。记余球边界上的同构为 \(\alpha:\pi^*E\to\pi^*F\)，则伴随的三元组为 \((\pi^*F,\pi^*E,\alpha^*)\)。

	原三元组的负元为 \((\pi^*F,\pi^*E,\alpha^{-1})\)。在边界上，\(\alpha^*\alpha\) 是 \(\pi^*E\) 的正定自同构，故
	\[
		\alpha_t=\bigl((1-t)I+t\alpha^*\alpha\bigr)\alpha^{-1},\qquad0\le t\le1,
	\]
	始终可逆，并连接 \(\alpha^{-1}\) 与 \(\alpha^*\)。所以两份三元组同伦，得到符号类的负号。

	分析一侧，闭像与椭圆正则性给出 \(\operatorname{coker}P\cong\ker P^*\) 以及 \(\operatorname{coker}P^*\cong\ker P\)。因而
	\[
		\operatorname{ind}_{\mathrm a}(P^*)
		=\dim\ker P^*-\dim\ker P
		=-\operatorname{ind}_{\mathrm a}(P).
	\]
	这里负号来自交换核与余核；拓扑一侧则来自交换三元组的两张丛并反转边界同构。指标定理将这两种说明联系起来。
\end{solution}
